\documentclass[11pt,a4paper]{article}
\usepackage{apnotes}
\setsunday{2}
\begin{document}

\notestitle{Deep Learning with PyTorch}{6--10}{Understand neurons, backpropagation, convolution and attention well enough to build them yourself; build a leakage-free image pipeline; test robustness honestly; and ship \textbf{P2 --- LeafDoctor}, a phone-sized plant-disease model.}

\section*{How to use these notes}
These notes follow the Sunday 2 slides topic by topic and explain each step in full, using the same colour code:
\legendbox
Numbers in \emph{Our notebook result} boxes come from \texttt{notebooks/S02\_Deep\_Learning\_PyTorch.ipynb}, run on a laptop CPU. The last section is your \textbf{Monday--Saturday project homework}.

\begin{tabularx}{\linewidth}{@{}c c X l@{}}
\toprule
\textbf{Hour} & \textbf{Mode} & \textbf{Focus} & \textbf{Dataset}\\
\midrule
6 & \teachtag & Neurons, autograd, backpropagation, training loop & \ds{Fashion-MNIST}\\
7 & \teachtag & Convolution, pooling, batch-norm, augmentation & \ds{CIFAR-10}\\
8 & \buildtag & Image pipeline, splits, leakage checks & \ds{PlantVillage (apple)}\\
9 & \buildtag & Vision Transformer, Grad-CAM, error analysis & \ds{PlantVillage (apple)}\\
10 & \buildtag & Distribution shift, quantisation, ONNX, model card & \ds{PlantVillage + shifts}\\
\bottomrule
\end{tabularx}

\begin{bdef}[Learning objectives]
By the end of today you can (1) explain neurons, activations, loss and backpropagation; (2) write a PyTorch training loop from scratch; (3) design a CNN and explain convolution, pooling and batch normalisation; (4) build a leakage-free image pipeline; (5) compare a CNN with a Vision Transformer and read Grad-CAM; (6) test robustness, quantise and export a model to ONNX.\\[2pt]
\textbf{Prerequisites:} Sunday 1 --- loss functions, gradient descent, train/validation/test splits.
\end{bdef}

%% =====================================================================
\hourhead{6}{TEACH}{Neural networks, autograd and backpropagation}{Fashion-MNIST (70,000 images)}{PyTorch core, autograd, custom training loops, loss curves}

\topic{The artificial neuron}
\begin{bdef}[Neural network]
Layers of simple units. Each unit takes a weighted sum of its inputs and passes it through a non-linear \emph{activation function}. Stacking layers lets the network represent very complex functions.
\end{bdef}
A single neuron is just logistic regression from Sunday 1. The power of neural networks comes from composing many of them: the first layer learns simple patterns, later layers combine those into more abstract ones. Because every operation is differentiable, the whole network can be trained with gradient descent.
\begin{bmath}[Neuron and activations]
$a=\varphi(w^\top x+b)$ \qquad ReLU $\varphi(z)=\max(0,z)$ \qquad Sigmoid $\sigma(z)=\frac{1}{1+e^{-z}}$\\[3pt]
Softmax turns a vector of scores into probabilities: $\;p_c=\dfrac{e^{z_c}}{\sum_{k}e^{z_k}}$.
\end{bmath}
\begin{bex}[One neuron by hand]
$x=(1,2)$, $w=(0.5,-1)$, $b=1$: $z=0.5-2+1=-0.5$. ReLU gives $0$; sigmoid gives $0.38$.
\end{bex}
\begin{bwarn}[Why activations are essential]
Without non-linear activations, a stack of linear layers is still one linear function ($W_2(W_1x)=(W_2W_1)x$), no matter how deep it is.
\end{bwarn}

\topic{Loss and the chain rule}
\begin{bmath}[Cross-entropy loss]
For the true class $y$ with predicted probability $p_y$: $\;L=-\log p_y$.
\end{bmath}
\begin{bex}[How harsh is cross-entropy?]
$p_y=0.9\Rightarrow L=0.105$; \;$p_y=0.1\Rightarrow L=2.30$. A confident wrong answer costs about 22 times more than a confident right one, which pushes the model hard away from over-confident mistakes.
\end{bex}
\begin{bdef}[Autograd]
PyTorch records every operation on tensors with \texttt{requires\_grad=True} in a \emph{computational graph}. Calling \texttt{backward()} walks the graph in reverse and applies the chain rule to compute every gradient automatically.
\end{bdef}
\begin{bmath}[Chain rule]
For $L(a(z(w)))$: $\;\dfrac{\partial L}{\partial w}=\dfrac{\partial L}{\partial a}\cdot\dfrac{\partial a}{\partial z}\cdot\dfrac{\partial z}{\partial w}$
\end{bmath}
\begin{bex}[Autograd check]
For $L=(2w-4)^2$ at $w=3$: $\frac{dL}{dw}=2(2w-4)\cdot2=8$. PyTorch returns exactly 8.
\end{bex}

\topic{Backpropagation and optimisers}
\begin{balgo}[Training with backpropagation]
\begin{algorithmic}[1]
\For{each epoch}
  \For{each mini-batch $(X_b,y_b)$}
    \State forward pass: $\hat y=f_\theta(X_b)$
    \State compute the loss $L(\hat y,y_b)$
    \State backward pass: $g=\nabla_\theta L$ via the chain rule
    \State update $\theta$ with the optimiser
  \EndFor
\EndFor
\end{algorithmic}
\end{balgo}
\begin{bmath}[Adam optimiser]
$m_t=\beta_1m_{t-1}+(1-\beta_1)g_t$ (moving average of gradients)\\
$v_t=\beta_2v_{t-1}+(1-\beta_2)g_t^2$ (moving average of squared gradients)\\
$\theta_t=\theta_{t-1}-\eta\,\hat m_t/(\sqrt{\hat v_t}+\epsilon)$, with bias-corrected $\hat m_t,\hat v_t$.
\end{bmath}
Adam gives each parameter its own effective step size: parameters with consistently large gradients take smaller steps, rarely-updated ones take larger steps. That is why it usually works with little tuning.
\begin{bdef}[Regularisation]
\textbf{Dropout} randomly zeroes units during training so the network cannot rely on any single one. \textbf{Weight decay} adds $\lambda\lVert\theta\rVert^2$ to the loss, keeping weights small.
\end{bdef}

\topic{The training loop in code}
\begin{bcode}
\begin{lstlisting}
mlp = nn.Sequential(nn.Linear(784, 256), nn.ReLU(), nn.Dropout(0.2),
                    nn.Linear(256, 128), nn.ReLU(), nn.Linear(128, 10))
opt = torch.optim.Adam(mlp.parameters(), lr=1e-3, weight_decay=1e-5)
for epoch in range(5):
    for xb, yb in batches:
        out  = mlp(xb)                       # 1. forward
        loss = F.cross_entropy(out, yb)      # 2. loss
        opt.zero_grad()                      # 3. clear old gradients
        loss.backward()                      # 4. backward
        opt.step()                           # 5. update
\end{lstlisting}
\end{bcode}
\begin{bres}
MLP (784--256--128--10) reaches \textbf{87.1\% test accuracy} on Fashion-MNIST after 5 epochs. Test loss falls from 0.464 to 0.357 while train loss falls from 0.660 to 0.332 --- no overfitting yet.
\end{bres}
\figw[0.85]{s2_mlp_curves.png}{Loss and accuracy curves. If train loss keeps falling while test loss rises, the model is memorising.}
\begin{bapp}
Image tagging in e-commerce (Fashion-MNIST was released by Zalando using its own catalogue images), handwriting recognition, and every deep-learning product built with PyTorch, TensorFlow or JAX.
\end{bapp}

%% =====================================================================
\hourhead{7}{TEACH}{Convolutional networks, augmentation, mixed precision}{CIFAR-10 (10,000-image split)}{CNN design, augmentation, AMP training}

\topic{Convolution}
\begin{bdef}[Convolutional layer]
A small filter (for example $3\times3$) slides over the image and computes a weighted sum at each position. The same weights are reused everywhere (\emph{weight sharing}) and each output sees only a small neighbourhood (\emph{locality}).
\end{bdef}
A fully connected layer on a $32\times32\times3$ image connects every pixel to every unit, ignoring that nearby pixels belong together. Convolution builds in two facts about images: patterns are local, and the same pattern (an edge, a spot) can appear anywhere. That makes CNNs far more data-efficient.
\begin{bmath}[Formulas]
$S(i,j)=\sum_m\sum_n I(i+m,\,j+n)\,K(m,n)+b$\\[3pt]
Output size $=\left\lfloor\dfrac{W-K+2P}{S}\right\rfloor+1$ \;(width $W$, kernel $K$, padding $P$, stride $S$)\\[3pt]
Parameters $=K\cdot K\cdot C_{in}\cdot C_{out}+C_{out}$
\end{bmath}
\begin{bex}[Worked example]
A $32\times32$ input with $K=3$, $P=1$, $S=1$ gives a $32\times32$ output. Going from 3 to 32 channels uses $3\cdot3\cdot3\cdot32+32=896$ parameters; a dense layer between the same sizes would need about 3 million.
\end{bex}

\topic{Pooling, batch normalisation, augmentation and mixed precision}
\begin{bdef}[Pooling]
Max-pooling keeps the largest value in each $2\times2$ window, halving height and width and making the network tolerant of small shifts.
\end{bdef}
\begin{bmath}[Batch normalisation]
$\hat x=\dfrac{x-\mu_B}{\sqrt{\sigma_B^2+\epsilon}}$, \;$y=\gamma\hat x+\beta$ --- normalise each channel over the mini-batch, then let the network rescale it with learned $\gamma,\beta$.
\end{bmath}
Batch normalisation keeps activations at a stable scale as earlier layers change, which allows higher learning rates and faster training.
\begin{balgo}[A CNN block stack]
\begin{algorithmic}[1]
\For{channels $c\in\{32,64,128\}$}
  \State Conv $3\times3$ $\rightarrow$ BatchNorm $\rightarrow$ ReLU
  \State Conv $3\times3$ $\rightarrow$ BatchNorm $\rightarrow$ ReLU
  \State MaxPool $2\times2$
\EndFor
\State Global average pooling $\rightarrow$ Dropout $\rightarrow$ Linear classifier
\end{algorithmic}
\end{balgo}
\begin{bdef}[Data augmentation and mixed precision]
\textbf{Augmentation} trains on randomly flipped and shifted copies, so the model learns the object rather than exact pixel positions. \textbf{Mixed precision} runs most maths in 16-bit (bfloat16), which is faster and lighter with almost the same predictions.
\end{bdef}
\begin{bcode}
\begin{lstlisting}
layers += [nn.Conv2d(cin, c, 3, padding=1), nn.BatchNorm2d(c), nn.ReLU(),
           nn.Conv2d(c, c, 3, padding=1), nn.BatchNorm2d(c), nn.ReLU(), nn.MaxPool2d(2)]
with torch.autocast(device_type="cpu", dtype=torch.bfloat16):   # mixed precision
    loss = F.cross_entropy(model(augment(xb)), yb)
\end{lstlisting}
\end{bcode}
\begin{bres}
A 289,194-parameter CNN trained for 8 epochs on 8,000 CIFAR-10 images reaches about \textbf{67\% accuracy} on 2,000 held-out images (peak 68.5\%). bfloat16 and fp32 give the same prediction on \textbf{99.0\%} of images.
\end{bres}
\begin{bwarn}[Normalise with training statistics only]
Compute per-channel mean and standard deviation on training images and apply them to validation and test. Using test statistics is a quiet form of leakage.
\end{bwarn}
\begin{bapp}
Photo search and visual search (Google Lens, Apple Photos), medical imaging such as diabetic-retinopathy and chest X-ray screening, factory visual inspection, and camera perception in driver-assistance systems.
\end{bapp}

%% =====================================================================
\hourhead{8}{BUILD}{LeafDoctor: image pipeline and leakage checks}{PlantVillage --- 4 apple classes, 3,171 images}{Datasets, DataLoaders, stratified splits, leakage checks}

\topic{Data pipelines and leakage}
\begin{bdef}[Data pipeline]
The steps that turn files into training batches: read, decode, resize, normalise, augment, batch. In PyTorch a \texttt{Dataset} returns one example and a \texttt{DataLoader} shuffles and batches them.
\end{bdef}
\begin{bdef}[Data leakage]
Information from the test set reaching training --- for example the same photo appearing in both. Scores then look excellent and mean nothing.
\end{bdef}
\begin{bmath}[Stratified split]
Each class keeps its share in every split: $\;\frac{n_c^{\,train}}{n^{\,train}}\approx\frac{n_c^{\,test}}{n^{\,test}}\approx\frac{n_c}{n}$. This matters here because \emph{cedar apple rust} has only 275 images.
\end{bmath}
\begin{balgo}[Leakage-safe split]
\begin{algorithmic}[1]
\For{each image}
  \State $h\gets$ MD5 hash of its pixels
\EndFor
\State keep only the first image of each hash
\State stratified split 70 / 15 / 15
\State \textbf{assert} no hash appears in both train and test
\end{algorithmic}
\end{balgo}
\figw[0.72]{s2_leaf_samples.png}{One example per class from the PlantVillage apple subset.}
\begin{bcode}
\begin{lstlisting}
hashes = [hashlib.md5(img.tobytes()).hexdigest() for img in images]
keep = first_index_per_unique(hashes)                      # de-duplicate BEFORE splitting
tr, tmp = train_test_split(keep, test_size=0.30, stratify=y[keep], random_state=42)
va, te  = train_test_split(tmp,  test_size=0.50, stratify=y[tmp],  random_state=42)
\end{lstlisting}
\end{bcode}
\begin{bres}
7 exact duplicates found. Before the fix, 1 duplicate had landed in both train and test. After de-duplicating first: 3,164 images, train/val/test = 2,214 / 475 / 475, \textbf{0 overlap}.
\end{bres}
\begin{bapp}
Crop-diagnosis apps such as Plantix and PlantNet and farm advisory services rely on exactly this kind of pipeline; the same pattern serves medical and retail image systems.
\end{bapp}
\begin{bwarn}[Background shortcut]
PlantVillage photos were taken on a plain lab background. A model can learn the setting instead of the disease --- keep this in mind when reading Grad-CAM and the robustness test.
\end{bwarn}

%% =====================================================================
\hourhead{9}{BUILD}{Vision Transformer, Grad-CAM and error analysis}{PlantVillage apple subset}{Vision Transformers, Grad-CAM, confusion matrices}

\topic{Vision Transformer (ViT)}
\begin{bdef}[Vision Transformer]
Cut the image into patches, treat each patch as a token, and process the sequence with a Transformer encoder. Self-attention lets every patch look at every other patch from the first layer.
\end{bdef}
\begin{bmath}[Patch tokens]
Number of patches $N=\frac{H}{P}\cdot\frac{W}{P}$; a $64\times64$ image with $P=8$ gives 64 tokens.\\[2pt]
Input sequence: $z_0=[x_{cls};\,x_1E;\,\dots;\,x_NE]+E_{pos}$ (patch embeddings plus position embeddings, with a learnable \texttt{[CLS]} token).
\end{bmath}
\begin{balgo}[ViT forward pass]
\begin{algorithmic}[1]
\State split the image into patches and embed each one
\State prepend the \texttt{[CLS]} token and add position embeddings
\State run $L$ encoder blocks (self-attention + MLP, with residuals)
\State classify from the final \texttt{[CLS]} vector
\end{algorithmic}
\end{balgo}
\begin{bcode}
\begin{lstlisting}
self.patch = nn.Conv2d(3, dim, kernel_size=8, stride=8)    # patch embedding
t = self.patch(x).flatten(2).transpose(1, 2)               # (B, 64, dim)
t = torch.cat([self.cls.expand(len(t), -1, -1), t], 1) + self.pos
return self.head(self.norm(self.encoder(t)[:, 0]))         # classify from [CLS]
\end{lstlisting}
\end{bcode}
\begin{bres}[Same data, same 8 epochs, held-out test set]
\begin{tabular}{lrr}
\toprule
Model & Parameters & Test accuracy\\
\midrule
Small CNN & 162,748 & \textbf{0.987}\\
Tiny ViT (from scratch) & 324,580 & 0.943\\
\bottomrule
\end{tabular}\\[3pt]
CNN per-class F1: apple scab 0.974, black rot 0.995, cedar apple rust 0.988, healthy 0.990.
\end{bres}
\begin{bdef}[Inductive bias]
The assumptions a model builds in. CNNs assume locality and translation invariance, which suits small image datasets. ViTs assume very little, so they need more data --- or \textbf{pretraining} --- to reach their potential.
\end{bdef}
\begin{bapp}
ViT-style encoders sit inside many modern image and multimodal models. In class on a GPU, fine-tune a pretrained model with \texttt{timm.create\_model("efficientnet\_b0", pretrained=True, num\_classes=4)} or a pretrained ViT from Hugging Face.
\end{bapp}

\topic{Grad-CAM: seeing what the model sees}
\begin{bdef}[Grad-CAM]
A heat map of the image regions that increased the score of the predicted class, computed from the gradients flowing into the last convolutional layer.
\end{bdef}
\begin{bmath}[Formula]
$\alpha_k=\frac{1}{Z}\sum_{i,j}\frac{\partial y^c}{\partial A^k_{ij}}$ (importance of feature map $k$) \qquad $L^c=\mathrm{ReLU}\big(\sum_k\alpha_kA^k\big)$
\end{bmath}
\begin{balgo}[Grad-CAM]
\begin{algorithmic}[1]
\State forward pass; keep the last convolutional feature maps $A^k$
\State backward pass from the predicted class score $y^c$
\State average the gradients over each map to get $\alpha_k$
\State weighted sum, ReLU, and upsample to image size
\end{algorithmic}
\end{balgo}
\figw[0.82]{s2_gradcam.png}{Grad-CAM on test leaves. For scab and black rot the heat sits on lesions; for healthy and cedar apple rust it spreads to the leaf edge and background --- a warning sign of shortcut learning.}
\begin{bwarn}[Accuracy can hide shortcuts]
A 98.7\% model can still be looking at the wrong thing. Grad-CAM is a quick way to check, especially in medical and safety-critical imaging.
\end{bwarn}

%% =====================================================================
\hourhead{10}{BUILD}{Robustness, quantisation, ONNX, model card}{PlantVillage + simulated field conditions}{Distribution shift, quantisation, ONNX, model cards}

\topic{Distribution shift and robustness}
\begin{bdef}[Distribution shift]
Deployment data differs from training data. Under \emph{covariate shift}, $p_{train}(x)\neq p_{field}(x)$ while the labelling rule stays the same --- for example darker, blurrier, noisier phone photos.
\end{bdef}
\begin{balgo}[Robustness test]
\begin{algorithmic}[1]
\For{each shift $s\in\{\text{dark},\text{blur},\text{noise},\text{rotate}\}$}
  \State $X_s\gets s(X_{test})$
  \State $\text{accuracy}_s\gets\text{evaluate}(f, X_s)$
\EndFor
\State report the worst case, not only the clean score
\end{algorithmic}
\end{balgo}
\begin{bres}[Test accuracy under simulated field conditions]
\begin{tabular}{lrr}
\toprule
Condition & CNN & Tiny ViT\\
\midrule
Clean (lab photos) & 0.987 & 0.943\\
Darker light ($\times0.6$) & 0.657 & 0.509\\
Blur ($3\times3$) & 0.851 & 0.952\\
Sensor noise ($\sigma=0.08$) & \textbf{0.200} & \textbf{0.918}\\
Rotated 90$^\circ$ & 0.958 & 0.916\\
\bottomrule
\end{tabular}
\end{bres}
\begin{bwarn}[The most important result of the day]
The ``best'' model on clean data collapses to 20\% --- worse than always answering \emph{healthy} --- with mild noise, while the ViT holds 92\%. Lab accuracy is not field accuracy. Fixes: train with noise, brightness and blur augmentation; collect real field photos; add a ``not sure'' answer for low confidence.
\end{bwarn}

\topic{Quantisation and ONNX}
\begin{bdef}[Quantisation]
Storing weights (and sometimes activations) as 8-bit integers instead of 32-bit floats: roughly $4\times$ smaller and faster on phones and CPUs.
\end{bdef}
\begin{bmath}[Affine quantisation]
$q=\mathrm{round}(x/s)+z$, \;$x\approx s\,(q-z)$, with scale $s$ and zero-point $z$.
\end{bmath}
\begin{bex}[Example]
Weights in $[-1,1]$ mapped to 256 levels: $s=2/255\approx0.0078$. With $z=0$, the weight $0.5$ becomes $q=\mathrm{round}(64.0)=64$.
\end{bex}
\begin{bdef}[ONNX]
Open Neural Network Exchange: a portable model format that runs in ONNX Runtime on servers, phones and browsers, independent of PyTorch.
\end{bdef}
\begin{bcode}
\begin{lstlisting}
q_vit = torch.ao.quantization.quantize_dynamic(vit, ff_layers, dtype=torch.qint8)
torch.onnx.export(cnn, example, "leafdoctor.onnx", input_names=["image"],
                  output_names=["logits"], dynamic_axes={"image": {0: "batch"}}, opset_version=17)
sess = onnxruntime.InferenceSession("leafdoctor.onnx")
\end{lstlisting}
\end{bcode}
\begin{bres}
Dynamic INT8 quantisation of the ViT: 1,286\,KB $\rightarrow$ 860\,KB, accuracy unchanged at 0.943. ONNX export of the CNN: 637\,KB, largest logit difference from PyTorch $2.4\times10^{-6}$, identical predictions.
\end{bres}

\topic{Model cards and on-device AI}
\begin{bdef}[Model card]
A short document that travels with a model: task, data, performance, robustness, known limits, intended use and uses that are out of scope.
\end{bdef}
\begin{bex}[LeafDoctor model card, in brief]
Test accuracy 0.987 on lab photos; 0.200--0.987 under simulated shifts; four apple classes only; offers ``not sure'' at low confidence; not a substitute for an agronomist; deployed as a 637\,KB ONNX file.
\end{bex}
\begin{bapp}
On-device vision such as offline translation and visual search on phones; edge cameras in farms and factories. Runtimes: ONNX Runtime, TensorFlow Lite, Core ML.
\end{bapp}

%% =====================================================================
\section*{Formula sheet}
\begin{tabularx}{\linewidth}{@{}l X@{}}
\toprule
Neuron / softmax & $a=\varphi(w^\top x+b)$, \; $p_c=e^{z_c}/\sum_k e^{z_k}$\\
Cross-entropy & $L=-\log p_y$\\
Adam & $\theta\leftarrow\theta-\eta\,\hat m_t/(\sqrt{\hat v_t}+\epsilon)$\\
Convolution & output $\lfloor (W-K+2P)/S\rfloor+1$; \; params $K^2C_{in}C_{out}+C_{out}$\\
Batch-norm & $y=\gamma\,(x-\mu_B)/\sqrt{\sigma_B^2+\epsilon}+\beta$\\
ViT tokens & $N=HW/P^2$\\
Grad-CAM & $L^c=\mathrm{ReLU}(\sum_k\alpha_kA^k)$\\
Quantisation & $q=\mathrm{round}(x/s)+z$\\
\bottomrule
\end{tabularx}

\section*{Check yourself}
\begin{enumerate}
\item What happens to a deep network if you remove every activation? \emph{It collapses into one linear model.}
\item A $64\times64$ RGB image passes a $3\times3$ conv with 16 filters, padding 1, stride 2. Output and parameters? \emph{$32\times32\times16$; $3\cdot3\cdot3\cdot16+16=448$.}
\item Why can one duplicated photo across train and test inflate accuracy? \emph{The model is scored on an image it memorised.}
\item Why did the ViT beat the CNN under noise but lose on clean images? \emph{Different inductive biases: global attention is less sensitive to pixel noise; the CNN fits local texture better on clean lab images.}
\item What does INT8 quantisation trade, and what did it cost here? \emph{Numeric precision for size and speed; accuracy stayed at 0.943.}
\end{enumerate}

\section*{Glossary}
\textbf{Epoch} one pass over the data $\cdot$ \textbf{Batch} examples per update $\cdot$ \textbf{Autograd} automatic gradients $\cdot$ \textbf{Kernel} convolution filter $\cdot$ \textbf{Augmentation} random training variations $\cdot$ \textbf{ViT} Vision Transformer $\cdot$ \textbf{Grad-CAM} class heat map $\cdot$ \textbf{Quantisation} low-precision weights $\cdot$ \textbf{ONNX} portable model format.

\begin{submitbox}
Deliverable \textbf{P2 --- LeafDoctor}: notebook, ONNX model (or download link), Grad-CAM gallery, robustness table and \texttt{MODEL\_CARD.md} in \texttt{P2-leafdoctor/}. Tag \texttt{sunday-2}, create a Release, submit the link.
\end{submitbox}

\clearpage
\section*{Monday--Saturday Project Homework: Field-Ready Plant Disease Classifier}
\addcontentsline{toc}{section}{Monday--Saturday Project Homework}
\begin{tcolorbox}[enhanced,colback=navy,colframe=gold,boxrule=0.8pt,arc=3pt,coltext=white,left=8pt,right=8pt,top=6pt,bottom=6pt]
{\color{gold}\bfseries Weekly project.} Lab photos are clean; farmers' photos are not. You will build a classifier that is honest about the gap between the two and closes as much of it as possible.\\[3pt]
{\color{gold}\bfseries Data:} PlantVillage (full, 38 classes) for training; PlantDoc (real field photos) for testing\\[3pt]
{\color{gold}\bfseries Rhythm:} about 2--3 hours a day. Commit every day; tag \texttt{week-2} and create a GitHub Release by \textbf{Saturday 23:59 local time}. Mentors review it before the next Sunday.
\end{tcolorbox}

\begin{tcolorbox}[enhanced,breakable,colback=white,colframe=sun,boxrule=0.8pt,arc=2pt,left=7pt,right=7pt,top=4pt,bottom=4pt,
  title={\bfseries Monday --- A 38-class image pipeline\hfill\normalfont\small 3 h},coltitle=white,colbacktitle=sun,fonttitle=\small]
\textbf{Goal.} Build a leakage-free data pipeline at full scale.
\begin{enumerate}[itemsep=1pt,topsep=2pt]
  \item Download PlantVillage (all 38 classes) and PlantDoc.
  \item De-duplicate by perceptual hash (not only exact MD5); report how many near-duplicates you found.
  \item Create stratified splits and a \texttt{torch.utils.data.Dataset} with augmentation (flip, rotate, colour jitter, random resized crop).
  \item Map PlantDoc labels onto PlantVillage classes; document classes with no match.
\end{enumerate}
\textbf{\color{dataC}Deliverable:} \texttt{src/data.py}, \texttt{reports/data\_card.md}\\
\textbf{\color{navy}Done when:} No hash shared across splits; class counts per split tabulated; augmentation visualised on 8 images.\\
\textbf{\color{skillC}Stretch:} Use \texttt{torchvision.transforms.v2} and benchmark loading speed with 0, 2 and 4 workers.
\end{tcolorbox}

\begin{tcolorbox}[enhanced,breakable,colback=white,colframe=sun,boxrule=0.8pt,arc=2pt,left=7pt,right=7pt,top=4pt,bottom=4pt,
  title={\bfseries Tuesday --- Transfer learning versus training from scratch\hfill\normalfont\small 3 h},coltitle=white,colbacktitle=sun,fonttitle=\small]
\textbf{Goal.} See what pretraining buys.
\begin{enumerate}[itemsep=1pt,topsep=2pt]
  \item On a GPU (Colab or Kaggle), train ResNet-18 from scratch for 10 epochs.
  \item Fine-tune an ImageNet-pretrained ResNet-18 for 10 epochs (freeze, then unfreeze).
  \item Plot learning curves; report accuracy and macro-F1 on the PlantVillage test split.
  \item Explain the difference in two paragraphs.
\end{enumerate}
\textbf{\color{dataC}Deliverable:} \texttt{notebooks/02\_transfer.ipynb} with curves and a comparison table\\
\textbf{\color{navy}Done when:} Same data, epochs and seed for both runs; macro-F1 reported.\\
\textbf{\color{skillC}Stretch:} Add a learning-rate range test and one-cycle schedule.
\end{tcolorbox}

\begin{tcolorbox}[enhanced,breakable,colback=white,colframe=sun,boxrule=0.8pt,arc=2pt,left=7pt,right=7pt,top=4pt,bottom=4pt,
  title={\bfseries Wednesday --- Modern backbones and mixed precision\hfill\normalfont\small 3 h},coltitle=white,colbacktitle=sun,fonttitle=\small]
\textbf{Goal.} Compare architectures fairly.
\begin{enumerate}[itemsep=1pt,topsep=2pt]
  \item Fine-tune EfficientNet-B0 and a small ViT with \texttt{timm}, using mixed precision.
  \item Record parameters, training time per epoch, GPU memory and test macro-F1 for each.
  \item Build a confusion matrix for the best model; list the 3 most confused class pairs.
\end{enumerate}
\textbf{\color{dataC}Deliverable:} Comparison table and confusion matrix in \texttt{reports/backbones.md}\\
\textbf{\color{navy}Done when:} All three models compared on identical splits; costs (time, memory) reported next to accuracy.\\
\textbf{\color{skillC}Stretch:} Try ConvNeXt-Tiny or a distilled ViT (DeiT).
\end{tcolorbox}

\begin{tcolorbox}[enhanced,breakable,colback=white,colframe=sun,boxrule=0.8pt,arc=2pt,left=7pt,right=7pt,top=4pt,bottom=4pt,
  title={\bfseries Thursday --- Robustness on real field photos\hfill\normalfont\small 3 h},coltitle=white,colbacktitle=sun,fonttitle=\small]
\textbf{Goal.} Measure and close the lab-to-field gap.
\begin{enumerate}[itemsep=1pt,topsep=2pt]
  \item Evaluate your best model on PlantDoc field photos; report the accuracy drop.
  \item Retrain with stronger augmentation (blur, noise, brightness, background cut-and-paste) and re-test.
  \item Build a Grad-CAM gallery of 12 field failures and group them by cause.
\end{enumerate}
\textbf{\color{dataC}Deliverable:} \texttt{reports/robustness.md} with before/after tables and the Grad-CAM gallery\\
\textbf{\color{navy}Done when:} Lab vs field accuracy reported for both training recipes; failure causes categorised.\\
\textbf{\color{skillC}Stretch:} Fine-tune on 10\% of PlantDoc and measure how much labelled field data helps.
\end{tcolorbox}

\begin{tcolorbox}[enhanced,breakable,colback=white,colframe=sun,boxrule=0.8pt,arc=2pt,left=7pt,right=7pt,top=4pt,bottom=4pt,
  title={\bfseries Friday --- Calibration, quantisation and edge deployment\hfill\normalfont\small 2--3 h},coltitle=white,colbacktitle=sun,fonttitle=\small]
\textbf{Goal.} Make it small, fast and honest about uncertainty.
\begin{enumerate}[itemsep=1pt,topsep=2pt]
  \item Measure calibration (expected calibration error) and apply temperature scaling.
  \item Add a 'not sure' answer below a confidence threshold; report coverage vs accuracy.
  \item Export to ONNX, quantise to INT8, and benchmark CPU latency (median and p95 over 200 images).
\end{enumerate}
\textbf{\color{dataC}Deliverable:} \texttt{models/leafdoctor\_int8.onnx}, \texttt{reports/deploy.md}\\
\textbf{\color{navy}Done when:} ECE before and after scaling; latency table; INT8 accuracy within 1 point of FP32 or explained.\\
\textbf{\color{skillC}Stretch:} Run the ONNX model in a browser with ONNX Runtime Web.
\end{tcolorbox}

\begin{tcolorbox}[enhanced,breakable,colback=white,colframe=sun,boxrule=0.8pt,arc=2pt,left=7pt,right=7pt,top=4pt,bottom=4pt,
  title={\bfseries Saturday --- Demo app and model card\hfill\normalfont\small 3 h},coltitle=white,colbacktitle=sun,fonttitle=\small]
\textbf{Goal.} Ship it.
\begin{enumerate}[itemsep=1pt,topsep=2pt]
  \item Build a Gradio or Streamlit app: upload a leaf photo, get the top-3 diseases with confidence and Grad-CAM.
  \item Write \texttt{MODEL\_CARD.md} covering lab vs field performance and the 'not sure' policy.
  \item Record a 5-minute demo; push, tag \texttt{week-2}, create a Release.
\end{enumerate}
\textbf{\color{dataC}Deliverable:} App, model card, demo link, Release \texttt{week-2}\\
\textbf{\color{navy}Done when:} App works on a phone photo; model card states field accuracy, not only lab accuracy.\\
\textbf{\color{skillC}Stretch:} Deploy on Hugging Face Spaces.
\end{tcolorbox}


\end{document}
